\section{An extension of the exponent-six exclusions}
\label{sec:genus-extension}

The earlier packet \texttt{grh\_transfers.zip} ends by asking whether
smaller subfields of a genus field, together with ramified ideals of
small norm, can improve its exponent-six exclusions. The following
argument answers that specific question. It uses the exported Dirichlet
GRH and an explicit prime bound; it does not extend the theta-lift
construction to a new family of $L$-functions. The mathematical deduction
is new relative to the five packets reviewed here. No publication
priority or Lean verification is claimed.

\subsection{A general certificate}

Let $D<-4$ be a negative fundamental discriminant, put $x=|D|$, and write
$K=\mathbb Q(\sqrt D)$ and $A=\operatorname{Cl}(K)$. A ramified ideal
will mean a squarefree product of the unique prime ideals above rational
primes dividing $D$. Let $t=\omega(x)$, the number of those rational primes.
The genus field $G$ satisfies
\[
 [G:\mathbb Q]=2^t,\qquad
 \operatorname{Gal}(G/K)\simeq A/A^2.
\]
It is multiquadratic over $\mathbb Q$ and unramified over $K$ at every
finite prime. These are the standard genus-field facts used in the
previous packet.

\begin{theorem}[Intermediate genus-field certificate]
\label{thm:genus-intermediate}
Suppose the exponent of $A$ divides $2m$, where $m$ is a positive odd
integer. Let $H\subseteq A[2]$ have dimension $k$ over $\mathbb F_2$.
Assume that each class in $H$ has a ramified ideal representative of norm
at most $R$. There is an abelian field $F/\mathbb Q$ containing $K$, of
degree $n=2^{t-k}$ and absolute discriminant $x^{n/2}$, such that every
unramified rational prime $p$ splitting completely in $F$ satisfies
\begin{equation}\label{eq:genus-norm}
                         x\le 4R p^m.
\end{equation}
Consequently the repository's Dirichlet GRH implies
\begin{equation}\label{eq:genus-universal}
 x\le 4R\bigl(2n\log x+\tfrac52n+5\bigr)^{2m}.
\end{equation}
Any sharper valid upper bound for the least completely split prime in
$F$ may replace the bound used in \eqref{eq:genus-universal}.
\end{theorem}

\begin{proof}
The map $A[2]\to A/A^2$ is an isomorphism. Indeed, $A^2$ has exponent
dividing the odd integer $m$, so it intersects $A[2]$ trivially. For
surjectivity, $a^m\in A[2]$ and $a^{1-m}$ is a square for every $a\in A$.
Regard $H$ as a subgroup of the genus Galois group, and put $F=G^H$.
Then $K\subseteq F$, $[F:\mathbb Q]=n$, and $F/K$ is unramified.
The discriminant tower formula gives
\[
                   |\Delta_F|=|D|^{[F:K]}=x^{n/2}.
\]
Moreover $F/\mathbb Q$ is abelian because $G/\mathbb Q$ is.

Let $p$ be as stated. Since $F$ contains $K$, the prime $p$ splits in
$K$ and is prime to $D$. Choose one prime $\mathfrak p$ of $K$ above $p$.
Its Artin symbol in $\operatorname{Gal}(G/K)$ lies in $H$, so
$[\mathfrak p]=h a^2$ for some $h\in H$ and $a\in A$. Raising to the
$m$th power gives $[\mathfrak p]^m=h$. Choose a ramified ideal $J$ in
class $h$ with $N(J)\le R$. Because $h^2=1$, the integral ideal
$\mathfrak p^mJ$ is principal, say $(\alpha)$.

The valuations of $(\alpha)$ at $\mathfrak p$ and
$\overline{\mathfrak p}$ are $m$ and zero: $J$ is supported only at
primes above $D$. Hence $\alpha\notin\mathbb Q$. Write
$\alpha=(a+b\sqrt D)/2$ with integers $a,b$ and $b\ne0$. Norms give
\[
 p^mN(J)=N_{K/\mathbb Q}(\alpha)
        =\frac{a^2+xb^2}{4}\ge\frac x4,
\]
proving \eqref{eq:genus-norm}. Notice that the norm cost is $R$, not
$R^m$: the ramified ideal is inserted \emph{after} taking the $m$th power.

The Dedekind zeta function of the abelian field $F$ factors into
Dirichlet $L$-functions, so the exported theorem supplies its RH.
Theorem 5.1 of Bach--Sorenson gives an unramified prime with identity
Frobenius and
\[
 p\le\bigl(4\log|\Delta_F|+\tfrac52n+5\bigr)^2.
\]
Substitution proves \eqref{eq:genus-universal}.
\end{proof}

The explicit prime theorem is in E. Bach and J. Sorenson,
\href{https://doi.org/10.1090/S0025-5718-96-00763-6}{\emph{Explicit bounds
for primes in residue classes}}, Math. Comp. 65 (1996), Theorem 5.1
and Table 3. Here it is used only for an abelian extension of
$\mathbb Q$, so no additional nonabelian GRH assumption is involved.

One may optimize the theorem by replacing $R$ with
\[
 R(H)=\max_{h\in H}\min\{N(J):J\text{ ramified and }[J]=h\},
\]
then varying $H$. The following explicit choice already gives a useful
improvement and avoids any search over subgroups.

\subsection{The rank issue, including the prime two}

Let $q_1<\cdots<q_t$ be the rational primes dividing $D$, with
ramified prime ideals $\mathfrak q_i$, and write
$P_j=q_1\cdots q_j$, with $P_0=1$.

\begin{lemma}\label{lem:ramified-independent}
Under the exponent hypothesis of Theorem~\ref{thm:genus-intermediate},
any selection of at most $t-2$ ramified prime classes is independent in
$A[2]$ and in $A/A^2$. Thus for $0\le k\le t-2$ the subgroup generated
by $[\mathfrak q_1],\ldots,[\mathfrak q_k]$ has dimension $k$ and has
ramified representatives of norm at most $P_k$.
\end{lemma}

\begin{proof}
Each ramified class has order dividing two. Suppose a nonempty
squarefree product $J$ of ramified ideals is principal, say $(\alpha)$.
Conjugation fixes $J$, so $\overline\alpha/\alpha$ is a unit of $K$.
Since $D<-4$, the units are $\{1,-1\}$. If
$\overline\alpha=\alpha$, then $\alpha$ is rational and its valuations
at all ramified primes are even, contrary to the definition of $J$.
Otherwise $\alpha$ is a rational multiple of $\sqrt D$.

Multiplication by a rational number changes every ramified valuation by
an even integer. At each odd prime dividing $D$, the valuation of
$\sqrt D$ is odd. At the prime above two its parity is the parity of
$v_2(D)$. A fundamental even discriminant has $v_2(D)=2$ or $3$.
Therefore a nonempty relation has support all $t$ ramified primes when
$D$ is odd or $v_2(D)=3$, and all $t-1$ odd ramified primes when
$v_2(D)=2$. No selection of at most $t-2$ can support a relation.
Finally, a class of order dividing two that lies in $A^2$ is trivial,
since $A^2$ has odd exponent. This proves independence also in the genus
quotient. Every class of the stated subgroup has a representative formed
from a subset of the first $k$ prime ideals, of norm at most $P_k$.
\end{proof}

\subsection{Four certificates in the unresolved exponent-six window}

\begin{theorem}[Additional necessary conditions in the census window]
\label{thm:genus-four}
Assume the repository's Dirichlet GRH. Let $D<-4$ be fundamental, let
$\operatorname{Cl}(\mathbb Q(\sqrt D))$ have exponent six, and suppose
\[
 3.1\cdot10^{20}<x=|D|\le2.5\cdot10^{25},\qquad t=\omega(x)\ge7.
\]
Then, with the products $P_j$ defined above,
\begin{equation}\label{eq:genus-four}
 P_{t-2}>1{,}200{,}000,\qquad
 P_{t-3}>75{,}000,\qquad
 P_{t-4}>850,\qquad
 P_{t-5}>14.
\end{equation}
\end{theorem}

\begin{proof}
Put $y=\log x$, so $47<y<59$. Apply
Theorem~\ref{thm:genus-intermediate} with $m=3$ and
$k=t-\log_2 n$, for $n=4,8,16,32$. The preceding lemma verifies the
rank in all four cases. Table 3 of Bach--Sorenson supplies
$\sqrt p\le M_n(y)$ with the following affine upper bounds:
\begin{center}
\begin{tabular}{rrrrr}
\toprule
$n$ & $k$ & $M_n(y)$ & Test value $R$ & Ratio at $y=47.2$\\
\midrule
4 & $t-2$ & $4.070y+6.38$ & $1{,}200{,}000$ & $<0.946745$\\
8 & $t-3$ & $6.516y+6.98$ & $75{,}000$ & $<0.937077$\\
16 & $t-4$ & $13.960y+6.44$ & $850$ & $<0.951538$\\
32 & $t-5$ & $27.920y+7.08$ & $14$ & $<0.977087$\\
\bottomrule
\end{tabular}
\end{center}
The last column bounds $4R M_n(47.2)^6/(3.1\cdot10^{20})$ and is
displayed only for readability; the decisive comparisons use rational
arithmetic.

For $n=8,16,32$, the logarithm of the field discriminant is
$(n/2)y$, lying respectively in $(188,236)$, $(376,472)$, and
$(752,944)$. The applicable Table-3 triples $(A,B,C)$ in
$\sqrt p\le A\log|\Delta_F|+Bn+C$ are
$(1.629,0.11,6.1)$ for $n=8$, and $(1.745,0.04,5.8)$ for $n=16,32$.
For $n=4$ and $2y\le100$, the triple is $(2.035,0.27,5.3)$.
For $2y\ge100$, the triple is $(1.527,0.17,6.4)$; its bound
$3.054y+7.08$ is below the displayed $4.070y+6.38$ when $y\ge50$.
Thus the stated envelopes cover the entire window.

For every row, $4R(ay+b)^6e^{-y}$ decreases when $y>6-b/a$, in
particular when $y>47$. Write $X_0=3.1\cdot10^{20}$. Since
$\log X_0<236/5$, the exact checks
\[
                        4R M_n(236/5)^6<X_0
\]
give $4R M_n(y)^6<x$ throughout $x>X_0$. If the corresponding $P_k$
were at most $R$, equation~\eqref{eq:genus-norm} would imply the reverse
inequality. This proves all four necessary conditions.
\end{proof}

The published EKN census and global bound are independent inputs to its
classification interpretation: Elsenhans--Kl\"uners--Nicolae,
\href{https://arxiv.org/pdf/1803.02056}{\emph{Imaginary quadratic number
fields with class groups of small exponent}}, Theorems 1 and 3,
Tables 1 and 2. They give 432 exponent-six fields through $X_0$ and,
under the quadratic RH supplied here, no exponent-six field above
$2.5\cdot10^{25}$. The preceding packet excludes the remaining cases
with $t\le6$ and those with $x/P^+(x)\le10^8$.
Theorem~\ref{thm:genus-four} therefore excludes additional families
from the remaining window. It does not enumerate all survivors or prove
completeness of the 432-field list without a family restriction.

\paragraph{A strict improvement over those stated conditions.}
The fundamental discriminant
\[
\begin{split}
D&=-8\cdot3\cdot5\cdot7\cdot11\cdot1000000007\cdot1000000009\\
 &=-9240000147840000582120
\end{split}
\]
has seven distinct ramified primes. It satisfies
$X_0<|D|<2.1\cdot10^{22}$ and
$|D|/P^+(|D|)=9240000064680>10^8$, so it passes the earlier packet's
stated numerical necessary conditions. But $P_5=2310<1{,}200{,}000$;
the new theorem rules out exponent six. This demonstrates a strict
enlargement of those exclusion families, not a claim that this individual
discriminant was previously difficult or unknown.

\paragraph{Verification.}
The accompanying \texttt{genus\_certificate.py} uses exact rational
arithmetic to verify all four endpoint comparisons and exponential-series
bounds for $47<\log X_0<47.2$ and $\log(2.5\cdot10^{25})<59$.
It certifies the two large example primes by trial division, and checks
fundamentality, the products and the prior numerical conditions.
The JSON output records the exact ratios. It does not replace the
mathematical proof or independently recheck the published census.
